\documentclass[12pt]{article}
\usepackage[T2A]{fontenc}
\usepackage[utf8]{inputenc}
\usepackage[russian]{babel}
\usepackage{amsmath}
\usepackage{amsfonts}
\usepackage{hyperref}
\usepackage[a4paper, margin=1in]{geometry}

\begin{document}

\section*{Пояснение к формуле ожидаемого результата}

\subsection*{Формула ожидаемого результата \(E\)}
Ожидаемый результат \(E\) определяется как:
\[
E = \frac{1}{1 + \exp\left(-g(\text{RD}_{\text{оппонента}})\frac{(r - r_{\text{оппонента}})}{400}\right)}
\]
где:
\begin{itemize}
    \item \(r\) — рейтинг игрока;
    \item \(r_{\text{оппонента}}\) — рейтинг оппонента;
    \item \(g(\text{RD}_{\text{оппонента}})\) — функция неопределенности, зависящая от рейтингового отклонения (RD) оппонента, задаваемая как:
    \[
    g(\text{RD}) = \frac{1}{\sqrt{1 + \frac{3\,\text{RD}^2}{\pi^2}}}
    \]
    \item \(\exp\) — экспоненциальная функция с основанием \(e\).
\end{itemize}

\subsection*{Обоснование выбора формулы}

\subsubsection*{Связь с системой Elo}
Формула \(E\) является модификацией классической логистической функции, используемой в системе Elo:
\[
E_{\text{Elo}} = \frac{1}{1 + 10^{\frac{(r_{\text{оппонента}} - r)}{400}}}
\]
В классической системе Elo разница рейтингов между игроками интерпретируется как вероятность победы одного над другим через логистическую зависимость. Коэффициент 400 масштабирует разницу рейтингов так, что 400 единиц соответствуют примерно десятикратному различию в ожидаемой силе игроков. Модифицированная версия заменяет основание 10 на экспоненту \(e\), что делает функцию более гибкой для интеграции с другими компонентами системы, такими как \(g(\text{RD})\).

\textbf{Источник}: \\
Elo, A. E. (1978). \textit{The rating of chessplayers, past and present}. \\
Доступно по адресу: \url{https://en.wikipedia.org/wiki/Elo_rating_system}.

\subsubsection*{Учет неопределенности через \(g(\text{RD})\)}
Введение функции \(g(\text{RD})\) заимствовано из системы Glicko-2. Эта функция уменьшает влияние ожидаемого результата, если рейтинг оппонента имеет высокую неопределенность (большое значение RD). Например, при \(\text{RD} \to 0\) (высокая точность рейтинга) \(g(\text{RD}) \to 1\), и формула приближается к стандартной логистической функции. При больших \(\text{RD}\) (высокая неопределенность) \(g(\text{RD}) < 1\), что снижает вес игры в обновлении рейтинга. Это делает систему устойчивой к чрезмерным колебаниям рейтинга у игроков с недостаточным количеством данных.

\textbf{Источник}: \\
Glickman, M. E. (2013). \textit{Example of the Glicko-2 system}. \\
Доступно по адресу: \url{http://glicko.net/glicko/glicko2.pdf}.

\subsubsection*{Математическая строгость и адаптивность}
Использование экспоненциальной формы и коэффициента \(g(\text{RD})\) позволяет интегрировать формулу с другими элементами системы Glicko-2, такими как обратная дисперсия (\(v^{-1}\)) и дельта (\(\Delta\)), которые вычисляются как:
\[
v^{-1} = g(\text{RD}_{\text{оппонента}})^2 \, E \, (1 - E)
\]
\[
\Delta = v \cdot g(\text{RD}_{\text{оппонента}}) \, (s - E)
\]
Здесь \(s\) — фактический результат игры (\(1\) — победа, \(0.5\) — ничья, \(0\) — поражение). Эти параметры обеспечивают адаптивное обновление рейтинга с учетом статистической значимости результата.

\textbf{Источник}: \\
Glickman, M. E. (2013). \textit{Example of the Glicko-2 system}. \\
Доступно по адресу: \url{http://glicko.net/glicko/glicko2.pdf}.

\subsection*{Преимущества формулы}
\begin{itemize}
    \item \textbf{Гибкость}: Учет неопределенности через \(g(\text{RD})\) позволяет системе адаптироваться к игрокам с разным уровнем данных.
    \item \textbf{Совместимость}: Формула легко интегрируется с классическими (Elo) и современными (Glicko-2) подходами.
    \item \textbf{Интерпретируемость}: Результат \(E\) остается вероятностью победы, что соответствует интуитивному пониманию рейтингов.
\end{itemize}

\subsection*{Заключение}
Формула ожидаемого результата \(E\) представляет собой продуманное сочетание традиционной логистической модели Elo и адаптивных элементов Glicko-2. Ее академическая обоснованность подтверждается использованием проверенных источников и математической строгостью, что делает ее подходящей для гибридной рейтинговой системы, обеспечивающей точность и устойчивость в оценке силы игроков.

\end{document}